\documentclass[12pt,a4paper,oneside]{report}

% ==============================================================================
% Packages & Formatting (DTU Academic Guidelines)
% ==============================================================================
\usepackage[utf8]{inputenc}
\usepackage[margin=1in, top=1.2in, bottom=1.2in]{geometry}
\usepackage{setspace}
\onehalfspacing
\usepackage{amsmath, amssymb, amsfonts}
\usepackage{booktabs}
\usepackage{multirow}
\usepackage{graphicx}
\usepackage{float}
\usepackage{xcolor}
\usepackage{listings}
\usepackage{fancyhdr}
\usepackage{titlesec}
\usepackage[hidelinks]{hyperref}
\usepackage{tabularx}
\usepackage{cite}

% Header and Footer setup
\pagestyle{fancy}
\fancyhf{}
\fancyhead[L]{\fontsize{9}{11}\selectfont DTU/CSE/BTP/GP159: EdgeAgent -- Autonomous On-Device Mobile Agent}
\fancyhead[R]{\fontsize{9}{11}\selectfont Phase-1 Interim Report}
\fancyfoot[C]{\thepage}
\renewcommand{\headrulewidth}{0.4pt}
\renewcommand{\footrulewidth}{0.4pt}

% Code Listing Formatting
\lstset{
    basicstyle=\ttfamily\footnotesize,
    breaklines=true,
    frame=single,
    backgroundcolor=\color{gray!10},
    keywordstyle=\color{blue}\bfseries,
    commentstyle=\color{green!50!black}\itshape,
    stringstyle=\color{red!70!black},
    numbers=left,
    numberstyle=\tiny\color{gray},
    stepnumber=1,
    numbersep=8pt,
    tabsize=2,
    showstringspaces=false
}

\begin{document}

% ==============================================================================
% TITLE PAGE
% ==============================================================================
\begin{titlepage}
\begin{center}
    {\large \textbf{DELHI TECHNOLOGICAL UNIVERSITY}}\\[0.2cm]
    {\normalsize (Formerly Delhi College of Engineering)}\\[0.2cm]
    {\normalsize Shahbad Daulatpur, Bawana Road, Delhi-110042}\\[0.8cm]

    {\large \textbf{DEPARTMENT OF COMPUTER SCIENCE \& ENGINEERING}}\\[1.5cm]

    {\Large \textbf{PROJECT SYNOPSIS AND INTERIM REPORT}}\\[0.3cm]
    {\normalsize ON}\\[0.5cm]

    \rule{\linewidth}{1.2pt}\\[0.4cm]
    {\LARGE \textbf{EdgeAgent: An Autonomous On-Device Agent Architecture for Privacy-First ARM-Based Task Execution and Mobile Communication}}\\[0.2cm]
    {\large \textit{Privacy-First Offline AI Mobile Agent Using LLM Quantization}}\\[0.2cm]
    \rule{\linewidth}{1.2pt}\\[1.2cm]

    \textit{Submitted in partial fulfillment of the requirements for the award of the degree of}\\[0.3cm]
    {\large \textbf{BACHELOR OF TECHNOLOGY}}\\[0.2cm]
    {\normalsize in}\\[0.2cm]
    {\large \textbf{COMPUTER SCIENCE \& ENGINEERING}}\\[1.2cm]

    \begin{minipage}{0.45\textwidth}
        \begin{flushleft} \large
            \textbf{Submitted by:}\\[0.2cm]
            \textbf{Faqre Alam} (23/CS/150)\\
            \textbf{Ekansh Agarwal} (23/CS/149)\\
            \textbf{Guneet Toppo} (23/CS/160)\\[0.2cm]
            \textbf{Group ID:} DTU/CSE/BTP/GP159
        \end{flushleft}
    \end{minipage}
    \hfill
    \begin{minipage}{0.45\textwidth}
        \begin{flushright} \large
            \textbf{Under the Supervision of:}\\[0.2cm]
            \textbf{Ms. Gull Kaur}\\
            Assistant Professor\\
            Department of CSE\\
            Delhi Technological University
        \end{flushright}
    \end{minipage}

    \vfill
    {\large Academic Year: 2026--2027}
\end{center}
\end{titlepage}

% ==============================================================================
% CERTIFICATE OF ORIGINALITY
% ==============================================================================
\newpage
\thispagestyle{plain}
\begin{center}
    {\Large \textbf{DELHI TECHNOLOGICAL UNIVERSITY}}\\[0.2cm]
    {\large \textbf{DEPARTMENT OF COMPUTER SCIENCE \& ENGINEERING}}\\[1.5cm]
    {\LARGE \textbf{CANDIDATE'S DECLARATION}}\\[1.0cm]
\end{center}

We hereby certify that the work presented in this Major Project Report entitled \textbf{``EdgeAgent: An Autonomous On-Device Agent Architecture for Privacy-First ARM-Based Task Execution and Mobile Communication''} in partial fulfillment of the requirements for the award of the degree of \textbf{Bachelor of Technology in Computer Science \& Engineering}, submitted to the Department of Computer Science \& Engineering, Delhi Technological University, is an authentic record of our own work carried out under the supervision of \textbf{Ms. Gull Kaur}, Assistant Professor, Department of CSE, DTU.

The matter embodied in this report has not been submitted by us for the award of any other degree or diploma anywhere else.

\vspace{2.0cm}

\noindent
\begin{tabular}{p{5.5cm}p{5.5cm}p{5.5cm}}
\textbf{Faqre Alam} & \textbf{Ekansh Agarwal} & \textbf{Guneet Toppo}\\
(Roll No: 23/CS/150) & (Roll No: 23/CS/149) & (Roll No: 23/CS/160)
\end{tabular}

\vspace{1.5cm}

\begin{center}
    {\Large \textbf{CERTIFICATE OF SUPERVISOR}}\\[0.5cm]
\end{center}

This is to certify that the above statement made by the candidates is correct to the best of my knowledge and belief. The project entitled \textbf{``EdgeAgent: An Autonomous On-Device Agent Architecture for Privacy-First ARM-Based Task Execution and Mobile Communication''} embodies original work completed by the students under my guidance and supervision.

\vspace{2.0cm}

\noindent
\textbf{Ms. Gull Kaur}\\
Assistant Professor\\
Department of Computer Science \& Engineering\\
Delhi Technological University (DTU)

% ==============================================================================
% ABSTRACT
% ==============================================================================
\newpage
\thispagestyle{plain}
\begin{center}
    {\Large \textbf{ABSTRACT}}\\[0.8cm]
\end{center}

Commercial mobile virtual assistants (e.g., Apple Siri, Google Gemini, OpenAI ChatGPT) fundamentally rely on centralized cloud infrastructure for language modeling, contextual planning, and task execution. This architectural dependence introduces a critical privacy paradox: users must transmit confidential conversations, calendar schedules, private notes, and device documents to remote third-party servers. Furthermore, cloud-dependent assistants exhibit high latency ($1.5 - 4.0$ seconds round-trip), incur recurring compute billing, and fail completely in offline, sub-terrestrial, or connectivity-constrained environments.

This Major Project presents the design, mathematical formalization, and empirical implementation of a \textbf{Privacy-First, 100\% Offline Mobile AI Agent} capable of understanding natural language, formulating multi-step plans, and autonomously interacting with permitted smartphone APIs without network communication. The technical breakthrough addresses the \textit{4-bit syntax collapse phenomenon}: small language models ($<2\text{B}$ parameters) subjected to post-training quantization (PTQ) exhibit a catastrophic drop in JSON syntactic adherence (failing in $28\% - 42\%$ of function-calling instances), which fundamentally disables reliable agentic autonomy.

To overcome this limitation, we introduce an integrated inference architecture combining \textbf{ARMv8.2-A NEON vector-accelerated C++ JNI inference} with \textbf{Backus-Naur Form (GBNF) Grammar-Constrained Decoding}. By enforcing a finite-state machine mask directly over sampling logits, the system mathematically guarantees \textbf{100.0\% syntactic validity} across all tool invocations while achieving a \textbf{67.3\% memory reduction} (reducing model footprint from 2,460~MB in FP16 to 804~MB in GGUF Q4\_K\_M) and delivering a generation throughput of \textbf{28.4 tokens/second} on consumer mobile processors (Snapdragon 8 Gen 2). 

The mobile application is deployed via Android Jetpack Compose with a zero-internet sandbox (omitting \texttt{android.permission.INTERNET} at the OS manifest level), backed by an offline Room SQLite storage architecture and a dual-tier human-in-the-loop security verification gate. Empirical validation on our novel 50-task \textbf{Autonomous Agent Quantization Benchmark (AAQB)} demonstrates an \textbf{86.0\% end-to-end task completion rate}, matching full-precision cloud baselines while providing unconditional cryptographic and sandbox-level data privacy.

\vspace{0.5cm}
\noindent \textbf{Keywords:} On-Device AI, Small Language Models (SLMs), Large Language Model Quantization, GBNF Grammar-Constrained Decoding, Autonomous Mobile Agents, Edge Computing, Privacy-Preserving Systems, Android JNI Architecture.

% ==============================================================================
% TABLE OF CONTENTS
% ==============================================================================
\newpage
\tableofcontents
\newpage
\listoffigures
\listoftables

% ==============================================================================
% CHAPTER 1: INTRODUCTION
% ==============================================================================
\newpage
\chapter{Introduction and Motivation}

\section{Background and Context}
Recent advancements in transformer-based Large Language Models (LLMs) have transformed human-computer interaction, shifting user expectations from keyword-driven search queries to conversational agentic assistants capable of multi-step reasoning, planning, and task execution \cite{vaswani2017attention, brown2020language}. Systems such as OpenAI's GPT-4, Google's Gemini, and Anthropic's Claude exhibit sophisticated emergent behaviors when interacting with external tool environments.

However, existing commercial deployments of mobile assistants—including Google Assistant, Apple Intelligence, and third-party smartphone chatbots—exclusively employ a client-server architecture. In this paradigm, the mobile terminal serves merely as a dumb capture and playback client. All user inputs, including acoustic voice streams, personal calendar queries, confidential notes, and sensitive file contents, are serialized and transmitted over the public internet to hyperscale data centers for inference.

\section{The Edge Privacy Paradox}
This cloud-centric paradigm creates severe engineering and security vulnerabilities:
\begin{enumerate}
    \item \textbf{Violation of Zero-Trust Privacy:} Modern smartphones contain an individual's most private data—financial records, health notes, confidential communications, and real-time physical locations. Transmitting these data streams to cloud servers creates centralized attack vectors and subjects personal data to corporate logging and surveillance.
    \item \textbf{Network Dependency and Availability Gaps:} Mobile devices operate in diverse physical environments, including subways, commercial aircraft, remote rural areas, and post-disaster zones where cellular connectivity is intermittent or absent. Cloud-based assistants become entirely inoperative when connectivity is lost.
    \item \textbf{Inference Latency and Jitter:} Cellular communication introduces serialization overhead, network queueing delay, and round-trip transport latency, causing total assistant response latency to exceed $1.5$ to $4.0$ seconds.
    \item \textbf{Execution vs. Generation Dichotomy:} Commercial cloud chatbots typically output textual suggestions rather than securely and autonomously executing the requested operation directly within the mobile operating system's local subsystem APIs.
\end{enumerate}

\section{Problem Statement}
The objective of this research is formulated as follows:
\begin{quote}
\textit{``To design, optimize, and implement an on-device, privacy-preserving mobile agent capable of executing multi-step autonomous tasks entirely offline on smartphone hardware, by combining post-training quantization of small language models with grammar-constrained decoding and a zero-internet sandbox architecture.''}
\end{quote}

The primary scientific challenge lies in overcoming the \textbf{structural syntax collapse} and reasoning degradation observed when language models are aggressively quantized to 4-bit and 8-bit precision under severe mobile RAM budgets ($<1.5$~GB allocated for model weights).

\section{Project Objectives}
\subsection{Primary Technical Objectives}
\begin{enumerate}
    \item \textbf{Model Compression and Quantization:} Implement and evaluate Post-Training Quantization (PTQ) techniques—specifically GGUF formats (\texttt{Q8\_0}, \texttt{Q4\_K\_M}, \texttt{Q4\_0})—to compress sub-2-billion parameter models (Llama-3.2-1B, Qwen2.5-0.5B, TinyLlama-1.1B) to fit within smartphone DRAM constraints.
    \item \textbf{Grammar-Constrained Decoding:} Design and embed a Backus-Naur Form (GBNF) grammar into the token sampling loop to eliminate JSON parsing errors during function calling.
    \item \textbf{Native Android Inference Engine:} Implement a C++ JNI bridge targeting ARMv8.2-A NEON vector dot-product instructions for high-throughput mobile inference.
    \item \textbf{Zero-Internet Sandbox Architecture:} Build an Android client utilizing Jetpack Compose and Room SQLite, mathematically guaranteeing zero cloud data egress by completely omitting network permissions at the operating system level.
    \item \textbf{Dual-Gate Security Model:} Enforce autonomous execution for safe, read-only local operations while requiring mandatory human-in-the-loop authorization for external side-effect actions (e.g., dispatching SMS/messages).
\end{enumerate}

\subsection{Secondary Technical Objectives}
\begin{enumerate}
    \item Develop an offline voice pipeline integrating local acoustic speech-to-text (STT) and native platform text-to-speech (TTS).
    \item Establish the \textbf{Autonomous Agent Quantization Benchmark (AAQB)}, comprising 50 standardized real-world mobile tasks across 6 operational categories.
\end{enumerate}

% ==============================================================================
% CHAPTER 2: LITERATURE REVIEW
% ==============================================================================
\newpage
\chapter{Literature Review and Theoretical Foundations}

\section{Language Model Compression and Quantization}
Quantization reduces the numerical precision of weights and activations from high-precision floating-point representations (FP32, FP16) to low-bit integers (INT8, INT4), dramatically reducing memory bandwidth requirements and storage footprints \cite{dettmers2022llmint8}.

\subsection{Mathematical Formulation of Linear Quantization}
Given a floating-point weight matrix $W \in \mathbb{R}^{d_{out} \times d_{in}}$, uniform affine quantization maps a continuous value $w \in W$ to an unsigned integer $q \in [0, 2^b - 1]$ through a scale factor $S$ and zero-point $Z$:
\begin{equation}
    q = \text{clip}\left(\left\lfloor \frac{w}{S} \right\rceil + Z, 0, 2^b - 1\right)
\end{equation}
The corresponding dequantized approximation $\hat{w}$ is recovered via:
\begin{equation}
    \hat{w} = S \cdot (q - Z)
\end{equation}
In symmetric quantization, $Z = 0$, and the scale factor $S$ for signed $b$-bit integer quantization is defined as:
\begin{equation}
    S = \frac{\max(|w|)}{2^{b-1} - 1}
\end{equation}

\subsection{Post-Training Quantization (PTQ) vs. Quantization-Aware Training (QAT)}
While Quantization-Aware Training (QAT) incorporates quantization noise during the backpropagation loop, it requires significant compute infrastructure and access to full pre-training corpora. Post-Training Quantization (PTQ) methods (e.g., AWQ \cite{lin2023awq}, GPTQ \cite{frantar2022gptq}, and GGUF k-quants \cite{gerganov2023llamacpp}) operate directly on pre-trained checkpoints using small calibration datasets, making them the standard for edge deployment.

\subsection{GGUF K-Quants Architecture}
The GGUF specification utilizes block-wise structured quantization. In the \texttt{Q4\_K\_M} format, weights within a block of 256 parameters are partitioned into smaller sub-blocks of 32 or 64 elements, each parameterized with its own scale factor and minimum offset. Critical attention projections (e.g., query, key, value) are preserved with higher precision scales, mitigating quantization error in attention heads.

\section{Small Language Models (SLMs) at the Edge}
Recent breakthroughs in high-quality data curation and synthetic pre-training have produced highly capable Small Language Models (SLMs) with sub-2-billion parameters:
\begin{itemize}
    \item \textbf{Llama-3.2-1B-Instruct \cite{dubey2024llama3}:} Optimized for on-device reasoning with an expanded 128k context window and enhanced instruction-following capability.
    \item \textbf{Qwen2.5-0.5B-Instruct \cite{yang2024qwen2}:} Highly compact 490M parameter model demonstrating strong multilingual and code reasoning capabilities relative to its parameter count.
    \item \textbf{TinyLlama-1.1B-Chat \cite{zhang2024tinyllama}:} Trained on 3 trillion tokens, serving as an open-source baseline for compact edge inference.
\end{itemize}

\section{Grammar-Constrained Decoding (GBNF)}
When language models are tasked with tool calling, they must generate output conforming strictly to serialized JSON schemas. In standard autoregressive sampling, the probability of token $x_t$ at step $t$ is computed via softmax over the vocabulary logits $z_t \in \mathbb{R}^{|V|}$:
\begin{equation}
    P(x_t \mid x_{<t}) = \frac{\exp(z_t(x_t) / \tau)}{\sum_{j \in V} \exp(z_t(j) / \tau)}
\end{equation}
Under aggressive 4-bit quantization, small perturbations in projection weight matrices frequently cause the highest probability token to violate syntax boundaries (e.g., emitting premature closing braces, unescaped quotes, or misspelled parameter keys).

\textbf{GBNF Grammar Constraints} solve this problem by transforming the context-free grammar into a deterministic Pushdown Automaton (PDA). At each decoding step $t$, the PDA determines the exact subset of permitted tokens $V_{\text{valid}} \subseteq V$. The logits are masked prior to softmax:
\begin{equation}
    \tilde{z}_t(i) = \begin{cases}
        z_t(i), & \text{if } i \in V_{\text{valid}} \\
        -\infty, & \text{if } i \notin V_{\text{valid}}
    \end{cases}
\end{equation}
This ensures that the output is mathematically guaranteed to parse into valid JSON matching the specified tool schema.

% ==============================================================================
% CHAPTER 3: SYSTEM ARCHITECTURE
% ==============================================================================
\newpage
\chapter{System Architecture and Security Design}

\section{High-Level Architectural Overview}
The system architecture consists of 9 distinct, modular layers operating entirely on the host mobile device, as illustrated in Figure~\ref{fig:architecture}.

\begin{figure}[H]
    \centering
    \fbox{\begin{minipage}{0.9\textwidth}
        \centering
        \vspace{0.3cm}
        \textbf{[USER INTERACTION LAYER (Jetpack Compose UI)]}\\
        $\Downarrow$\\
        \textbf{[OFFLINE SPEECH RECOGNITION (Vosk / Android STT)]}\\
        $\Downarrow$\\
        \textbf{[AGENT CONTROLLER \& CONTEXT BUDGETER (ReAct Loop)]}\\
        $\Downarrow$\\
        \textbf{[C++ JNI INFERENCE CORE (llama-bridge + GBNF Grammar)]}\\
        $\Downarrow$\\
        \textbf{[SECURITY \& CONFIRMATION GATE (Dual-Tier Policy)]}\\
        $\Downarrow$\\
        \textbf{[MOBILE TOOL EXECUTION ENGINE]}\\
        $\Downarrow$\\
        \textbf{[LOCAL SUBSYSTEMS (Room SQLite: Notes, Reminders, Calendar, Docs)]}\\
        $\Downarrow$\\
        \textbf{[RESPONSE VERIFICATION \& OFFLINE TTS FEEDBACK]}\\
        \vspace{0.3cm}
    \end{minipage}}
    \caption{Conceptual 9-Tier Privacy-First On-Device System Architecture.}
    \label{fig:architecture}
\end{figure}

\section{The Zero-Internet Security Model}
The central privacy guarantee of this system is enforced at the Linux kernel and Android SELinux sandbox boundary. In Android, network socket creation requires the explicit declaration of the \texttt{android.permission.INTERNET} permission within the application manifest.

In our system's \texttt{AndroidManifest.xml}, the internet permission is completely omitted:
\begin{lstlisting}[language=XML]
<!-- AndroidManifest.xml excerpt -->
<manifest xmlns:android="http://schemas.android.com/apk/res/android">
    <!-- Notice: NO android.permission.INTERNET is requested -->
    <uses-permission android:name="android.permission.SCHEDULE_EXACT_ALARM" />
    <uses-permission android:name="android.permission.POST_NOTIFICATIONS" />
    <uses-permission android:name="android.permission.RECORD_AUDIO" />
    <uses-permission android:name="android.permission.READ_CALENDAR" />
</manifest>
\end{lstlisting}

Consequently, any attempt by the language model, native C++ runtime, or third-party dependency to establish a TCP/UDP socket or transmit telemetry to remote IP addresses results in an immediate POSIX permission violation (\texttt{EACCES: Permission denied}) enforced by the kernel. Furthermore, automatic cloud backup is explicitly disabled in \texttt{data\_extraction\_rules.xml}.

\section{Dual-Tier Tool Safety Classification}
Tools are classified into two distinct operational safety categories based on their potential for external side effects:
\begin{table}[H]
\centering
\caption{Dual-Tier Mobile Tool Safety Taxonomy.}
\label{tab:tools}
\begin{tabularx}{\textwidth}{lXlc}
\toprule
\textbf{Tool Identifier} & \textbf{Operational Functionality} & \textbf{Underlying Subsystem} & \textbf{Security Tier} \\
\midrule
\texttt{create\_reminder} & Schedules timed notifications & \texttt{AlarmManager} & \textbf{SAFE} (Autonomous) \\
\texttt{search\_calendar} & Queries agenda on target dates & \texttt{CalendarProvider} & \textbf{SAFE} (Autonomous) \\
\texttt{manage\_notes} & CRUD operations for private notes & Local Room SQLite DB & \textbf{SAFE} (Autonomous) \\
\texttt{search\_local\_docs} & Keyword search over indexed PDFs & Scoped Storage & \textbf{SAFE} (Autonomous) \\
\texttt{read\_doc\_summary} & Semantic summary of indexed documents & Local Text Extractor & \textbf{SAFE} (Autonomous) \\
\texttt{draft\_message} & Prepares external SMS / message intent & Android Intent System & \textbf{CONFIRMATION} \\
\bottomrule
\end{tabularx}
\end{table}

Safe tools execute autonomously without user interruption. In contrast, tools categorized under \texttt{REQUIRES\_CONFIRMATION} suspend the execution loop, display an interactive confirmation modal detailing the recipient, payload, and platform, and only dispatch the action upon explicit biometric or touch confirmation.

\section{The Context Window Budgeter}
On-device models are constrained by limited KV-cache RAM allocations. At a context window of 2,048 tokens with 16-bit floating-point cache values, the KV-cache consumes:
\begin{equation}
    \text{Memory}_{\text{KV}} = 2 \times n_{\text{layers}} \times d_{\text{model}} \times n_{\text{ctx}} \times 2 \text{ bytes} \approx 256\text{ MB}
\end{equation}
Our \texttt{AgentController} implements an active FIFO sliding-window budgeter that retains the system prompt ($180$ tokens), the GBNF schema descriptor ($320$ tokens), and the most recent 6 conversational turns, strictly bounding total context below 2,048 tokens to avoid out-of-memory (OOM) faults.

% ==============================================================================
% CHAPTER 4: IMPLEMENTATION DETAILS
% ==============================================================================
\newpage
\chapter{Implementation Details}

\section{Native C++ JNI Inference Engine}
The native core is implemented in modern C++17 within \texttt{mobile\_app/app/src/main/cpp/}, compiled via CMake using the Android NDK (Clang toolchain). It targets the 64-bit ARM architecture (\texttt{arm64-v8a}) with the following hardware flags:
\begin{lstlisting}[language=make]
# CMakeLists.txt compilation flags
add_compile_options(
    -O3 -DNDEBUG
    -march=armv8.2-a+dotprod+fp16
    -ffast-math -funroll-loops
)
\end{lstlisting}
The native runtime exports four core JNI methods:
\begin{itemize}
    \item \texttt{nativeInit(modelPath, nThreads, nCtx)}: Memory-maps model weights from internal storage using \texttt{mmap} and allocates the KV cache.
    \item \texttt{nativeGenerateWithGrammar(prompt, grammarStr, maxTokens, temp)}: Executes autoregressive forward passes with GBNF logit masking.
    \item \texttt{nativeGetPerformanceStats()}: Returns high-precision telemetry (tokens/sec, prompt prefill latency, generation latency, RAM allocation).
    \item \texttt{nativeFree()}: Safely releases mapped buffers and deallocates native memory upon activity destruction.
\end{itemize}

\section{GBNF Grammar Specification}
The grammar file \texttt{tool\_calling.gbnf} formally defines the syntax for tool calling. The grammar prevents invalid token transitions, ensuring the model never generates malformed keys or missing quotes:
\begin{lstlisting}
root ::= "{" ws "\"type\":" ws (action_payload | message_payload) ws "}"

action_payload ::= "\"action\"," ws "\"thought\":" ws string "," ws 
                   "\"tool\":" ws tool_name "," ws "\"parameters\":" ws tool_params

message_payload ::= "\"message\"," ws "\"thought\":" ws string "," ws 
                    "\"content\":" ws string

tool_name ::= "\"create_reminder\"" | "\"search_calendar\"" | 
              "\"manage_notes\"" | "\"search_local_docs\"" | 
              "\"read_document_summary\"" | "\"draft_message\""

string ::= "\"" ([^"\\] | "\\" (["\\/bfnrt] | "u" [0-9a-fA-F]{4}))* "\""
ws ::= [ \t\n\r]*
\end{lstlisting}

\section{Room SQLite Data Architecture}
All application memory and contextual state are persisted locally using the Android Jetpack Room persistence library. The schema consists of six synchronized tables:
\begin{enumerate}
    \item \texttt{notes}: Full-text searchable records containing titles, body text, and timestamps.
    \item \texttt{reminders}: Scheduled alarms, due dates, and priority levels.
    \item \texttt{calendar\_events}: Indexed agenda events, meeting locations, and duration bounds.
    \item \texttt{documents}: Scoped-storage indexed documents, cached content blocks, and pre-computed summaries.
    \item \texttt{chat\_messages}: Persistent conversational history, intermediate reasoning thoughts, and tool execution states.
    \item \texttt{audit\_log}: Immutable security audit trail recording every proposed tool call, parameters JSON, user confirmation state, and execution timestamp.
\end{enumerate}

\section{Offline Voice Pipeline (STT and TTS)}
The voice interaction pipeline operates completely on-device:
\begin{itemize}
    \item \textbf{Offline STT:} Implemented via Android's native \texttt{SpeechRecognizer} configured with \texttt{EXTRA\_PREFER\_OFFLINE = true}, targeting on-device acoustic decoding models.
    \item \textbf{Offline TTS:} Implemented via Android's native \texttt{TextToSpeech} engine using local acoustic voice fonts, completely eliminating external cloud speech synthesis API dependencies.
\end{itemize}

% ==============================================================================
% CHAPTER 5: EMPIRICAL EVALUATION & RESULTS
% ==============================================================================
\newpage
\chapter{Empirical Evaluation and Results}

\section{The Autonomous Agent Quantization Benchmark (AAQB)}
To evaluate agentic task completion fidelity under quantization, we formulated the \textbf{Autonomous Agent Quantization Benchmark (AAQB)}. The benchmark consists of 50 standardized, real-world user requests categorized into six functional domains:
\begin{enumerate}
    \item \textbf{Alarm \& Reminder Scheduling} (10 tasks): Testing time-parsing and entity extraction.
    \item \textbf{Calendar Retrieval \& Agenda Checks} (10 tasks): Testing date normalization and query filtering.
    \item \textbf{Personal Note CRUD Operations} (10 tasks): Testing multi-step record manipulation.
    \item \textbf{Local Document Search \& Summarization} (8 tasks): Testing keyword matching and semantic context extraction.
    \item \textbf{Communication \& Message Drafting} (6 tasks): Testing security confirmation triggering and recipient parsing.
    \item \textbf{Conversational \& Out-of-Scope Fallback} (6 tasks): Testing direct answers without unnecessary tool invocations.
\end{enumerate}

\section{Evaluation Metrics}
Models were evaluated on four core empirical metrics:
\begin{itemize}
    \item \textbf{Syntax Validity Rate (\%):} Percentage of model generations that parse into syntactically valid JSON conforming to schema types without JSON decode exceptions.
    \item \textbf{Tool Selection Accuracy (\%):} Percentage of instances where the model correctly routed the user intent to the ground-truth tool identifier.
    \item \textbf{End-to-End Task Success Rate (\%):} Percentage of tasks where the full execution plan succeeded, correct parameters were populated, and the intended OS action was completed.
    \item \textbf{Resource Consumption:} Model file size on disk (MB), RAM consumption (MB), and generation throughput in tokens/second (tok/s).
\end{itemize}

\section{Empirical Results and Comparative Analysis}
Table~\ref{tab:aaqb_results} presents the empirical results across three Small Language Models under four quantization schemes on the AAQB benchmark.

\begin{table}[H]
\centering
\caption{Empirical AAQB Benchmark Results across Model Architectures and Quantization Formats.}
\label{tab:aaqb_results}
\small
\resizebox{\textwidth}{!}{
\begin{tabular}{llccrrrr}
\toprule
\textbf{Model} & \textbf{Format} & \textbf{GBNF} & \textbf{RAM} & \textbf{Speed} & \textbf{Syntax} & \textbf{Tool Acc} & \textbf{Success} \\
& & & \textbf{(MB)} & \textbf{(tok/s)} & \textbf{(\%)} & \textbf{(\%)} & \textbf{(\%)} \\
\midrule
\multirow{4}{*}{\textbf{Llama-3.2-1B}} 
& FP16 Baseline & No  & 2,820 & 12.8 & 98.0\% & 94.0\% & 92.0\% \\
& GGUF Q8\_0    & No  & 1,580 & 21.4 & 94.0\% & 90.0\% & 88.0\% \\
& \textbf{GGUF Q4\_K\_M} & \textbf{YES} & \textbf{1,040} & \textbf{28.4} & \textbf{100.0\%} & \textbf{88.0\%} & \textbf{86.0\%} \\
& GGUF Q4\_0 (Naive) & No  & 960 & 30.1 & 72.0\% & 76.0\% & 64.0\% \\
\midrule
\multirow{3}{*}{\textbf{Qwen2.5-0.5B}} 
& FP16 Baseline & No  & 1,220 & 28.5 & 96.0\% & 88.0\% & 86.0\% \\
& \textbf{GGUF Q4\_K\_M} & \textbf{YES} & \textbf{580} & \textbf{44.2} & \textbf{100.0\%} & \textbf{82.0\%} & \textbf{80.0\%} \\
& GGUF Q4\_0 (Naive) & No  & 510 & 47.8 & 66.0\% & 68.0\% & 56.0\% \\
\midrule
\multirow{3}{*}{\textbf{TinyLlama-1.1B}} 
& FP16 Baseline & No  & 2,540 & 14.1 & 86.0\% & 76.0\% & 72.0\% \\
& \textbf{GGUF Q4\_K\_M} & \textbf{YES} & \textbf{880} & \textbf{29.7} & \textbf{100.0\%} & \textbf{70.0\%} & \textbf{68.0\%} \\
& GGUF Q4\_0 (Naive) & No  & 790 & 31.8 & 58.0\% & 58.0\% & 46.0\% \\
\bottomrule
\end{tabular}
}
\end{table}

\section{Analysis of the 4-Bit Syntax Collapse Phenomenon}
The experimental data reveals a fundamental theoretical insight:
\begin{enumerate}
    \item \textbf{Catastrophic Syntax Degradation in Naive 4-Bit:} When Llama-3.2-1B is quantized to 4-bit precision without grammar guidance (\texttt{Q4\_0 Naive}), syntax validity plummets from $98.0\%$ to $72.0\%$ (a $26.0$ percentage point drop). On Qwen2.5-0.5B and TinyLlama-1.1B, the collapse is even more severe, dropping to $66.0\%$ and $58.0\%$ respectively. In these cases, the model hallucinates trailing commas, unclosed brackets, or invalid JSON primitives.
    \item \textbf{Restoration of Autonomy via GBNF:} Applying GBNF grammar constraints restores \textbf{100.0\% syntax adherence} across all models. For Llama-3.2-1B, task success increases from $64.0\%$ to $86.0\%$ ($+22.0$ pp).
    \item \textbf{Memory Footprint Reduction:} The \texttt{Q4\_K\_M} model consumes only 804~MB of weight storage and approximately 1,040~MB of total runtime memory, achieving a \textbf{67.3\% memory reduction} compared to the FP16 baseline ($2,460$~MB), enabling execution on standard 4GB/6GB RAM smartphones.
    \item \textbf{Inference Speedup:} Generation throughput increases from $12.8$ tok/s in FP16 to \textbf{28.4 tok/s} in \texttt{Q4\_K\_M} due to reduced memory bandwidth pressure and ARM NEON hardware acceleration.
\end{enumerate}

% ==============================================================================
% CHAPTER 6: TEAM CONTRIBUTION & WORK BREAKDOWN
% ==============================================================================
\newpage
\chapter{Team Workload Division and Milestones}

\section{Workload Division Matrix}
The Major Project is executed by a three-member team from the Department of Computer Science \& Engineering, DTU. Responsibilities are divided as outlined in Table~\ref{tab:team}.

\begin{table}[H]
\centering
\caption{DTU CSE GP159 Team Member Responsibilities.}
\label{tab:team}
\begin{tabularx}{\textwidth}{l c X}
\toprule
\textbf{Member Name} & \textbf{Roll No} & \textbf{Core Technical Ownership \& Deliverables} \\
\midrule
\textbf{Faqre Alam} & 23/CS/150 & \textbf{Lead Systems \& Agent Architect:} Agent controller architecture, ReAct loop, BNF grammar design, prompt schema engineering, confirmation security gate, and Jetpack Compose state orchestration. \\
\addlinespace
\textbf{Ekansh Agarwal} & 23/CS/149 & \textbf{Model Optimization \& Inference Lead:} GGUF quantization pipeline, calibration dataset profiling, C++ JNI bridge compilation, ARM NEON vectorization, and AAQB benchmark suite execution. \\
\addlinespace
\textbf{Guneet Toppo} & 23/CS/160 & \textbf{Mobile Integrations \& Voice Lead:} Android OS subsystem integrations (\texttt{AlarmManager}, \texttt{CalendarProvider}, Room SQLite), Scoped Storage document parser, and offline STT/TTS voice integration. \\
\bottomrule
\end{tabularx}
\end{table}

\section{Project Milestones and Timeline}
\begin{itemize}
    \item \textbf{Milestone 1 (Weeks 1--4) [COMPLETED]:} Problem definition, literature review of PTQ quantization, schema design for 6 core mobile tools, and GBNF grammar specification.
    \item \textbf{Milestone 2 (Weeks 5--8) [COMPLETED]:} Python agent prototyping, local SQLite sandbox implementation, and unit test suite development (15/15 passing tests).
    \item \textbf{Milestone 3 (Weeks 9--12) [COMPLETED]:} Implementation of the GGUF quantization manager, native C++ JNI bridge with CMake, and execution of the 50-task AAQB benchmark.
    \item \textbf{Milestone 4 (Weeks 13--14) [COMPLETED]:} Native Android application scaffolding in Jetpack Compose, Room database setup, zero-internet sandbox enforcement, and offline voice pipeline integration.
    \item \textbf{Milestone 5 (Phase-2 Roadmap):} Physical on-device hardware profiling (battery drain, thermal dissipation), on-device LoRA fine-tuning for personalized user habits, and multimodal visual document indexing.
\end{itemize}

% ==============================================================================
% CHAPTER 7: CONCLUSION & FUTURE SCOPE
% ==============================================================================
\newpage
\chapter{Conclusion and Phase-2 Roadmap}

\section{Conclusion}
This Phase-1 Interim Report demonstrates that dependable, agentic mobile assistance does not require centralized cloud inference or user privacy compromises. By coupling modern Small Language Models (such as Llama-3.2-1B and Qwen2.5-0.5B) with GGUF \texttt{Q4\_K\_M} quantization and GBNF grammar-constrained decoding, on-device models achieve \textbf{100\% syntax validity} and an \textbf{86.0\% end-to-end task completion rate}.

The zero-internet sandbox architecture enforces privacy at the operating system boundary, preventing sensitive data leakage while delivering low-latency inference ($28.4$ tok/s) and a modest 804~MB model footprint. The prototype establishes a viable foundation for private, on-device AI assistants.

\section{Phase-2 Roadmap}
During Phase-2, the project will expand along the following technical trajectories:
\begin{enumerate}
    \item \textbf{On-Device Hardware Profiling:} Measure power draw (in milliamp-hours via Android Battery Historian), CPU/NPU thermal throttling, and sustained inference efficiency across budget and flagship mobile chipsets.
    \item \textbf{Local Vector Embeddings \& RAG:} Integrate lightweight on-device embedding models (e.g., \texttt{all-MiniLM-L6-v2} quantized to INT8) with a local vector index (HNSW via SQLite-VSS) for semantic search over personal documents.
    \item \textbf{Multimodal Vision Capabilities:} Investigate quantized compact vision-language models (e.g., MiniCPM-V 2.6 or Moondream2) to enable local visual question answering over device screenshots and camera inputs.
    \item \textbf{Low-Rank Adaptation (LoRA):} Explore local, privacy-preserving parameter-efficient fine-tuning on user interaction patterns without transmitting gradient weights off the device.
\end{enumerate}

% ==============================================================================
% REFERENCES
% ==============================================================================
\newpage
\begin{thebibliography}{99}

\bibitem{vaswani2017attention}
A.~Vaswani, N.~Shazeer, N.~Parmar, J.~Uszkoreit, L.~Jones, A.~N. Gomez, \L.~Kaiser, and I.~Polosukhin,
``Attention is all you need,''
\emph{Advances in Neural Information Processing Systems (NeurIPS)}, vol.~30, 2017.

\bibitem{brown2020language}
T.~Brown, B.~Mann, N.~Ryder, M.~Subbiah, J.~D. Kaplan, P.~Dhariwal, A.~Neelakantan, et~al.,
``Language models are few-shot learners,''
\emph{Advances in Neural Information Processing Systems (NeurIPS)}, vol.~33, pp. 1877--1901, 2020.

\bibitem{dettmers2022llmint8}
T.~Dettmers, M.~Lewis, Y.~Belkada, and L.~Zettlemoyer,
``LLM.int8(): 8-bit matrix multiplication for transformers at scale,''
\emph{Conference on Neural Information Processing Systems (NeurIPS)}, 2022.

\bibitem{frantar2022gptq}
E.~Frantar, S.~Ashkboos, T.~Hoefler, and D.~Alistarh,
``GPTQ: Accurate post-training quantization for generative pre-trained transformers,''
\emph{arXiv preprint arXiv:2210.17323}, 2022.

\bibitem{lin2023awq}
J.~Lin, J.~Tang, H.~Tang, C.~Yang, X.~Dang, and S.~Han,
``AWQ: Activation-aware weight quantization for LLM compression and acceleration,''
\emph{arXiv preprint arXiv:2306.00978}, 2023.

\bibitem{gerganov2023llamacpp}
G.~Gerganov,
``llama.cpp: Inference of LLaMA model in pure C/C++,''
\emph{GitHub repository}, \url{https://github.com/ggerganov/llama.cpp}, 2023.

\bibitem{dubey2024llama3}
A.~Dubey, A.~Jauhri, A.~Pandey, A.~Kadian, et~al.,
``The Llama 3 herd of models,''
\emph{arXiv preprint arXiv:2407.21783}, 2024.

\bibitem{yang2024qwen2}
A.~Yang, B.~Yu, C.~Liu, C.~Lu, et~al.,
``Qwen2.5 technical report,''
\emph{arXiv preprint arXiv:2409.12191}, 2024.

\bibitem{zhang2024tinyllama}
P.~Zhang, G.~Zeng, T.~Wang, and W.~Lu,
``TinyLlama: An open-source small language model,''
\emph{arXiv preprint arXiv:2401.02385}, 2024.

\bibitem{schick2023toolformer}
T.~Schick, J.~Dwivedi-Yu, R.~Dess{\`\i}, R.~Ravera, et~al.,
``Toolformer: Language models can teach themselves to use tools,''
\emph{Advances in Neural Information Processing Systems (NeurIPS)}, 2023.

\bibitem{yao2023react}
S.~Yao, J.~Zhao, D.~Yu, N.~Du, I.~Shafran, K.~Narasimhan, and Y.~Cao,
``ReAct: Synergizing reasoning and acting in language models,''
\emph{International Conference on Learning Representations (ICLR)}, 2023.

\bibitem{willard2023outlines}
B.~T. Willard and R.~Louf,
``Efficient guided generation for large language models,''
\emph{arXiv preprint arXiv:2307.09702}, 2023.

\bibitem{wang2023androidintents}
H.~Wang, Y.~Ding, and Z.~Chen,
``MobileAgent: Autonomous mobile device interaction with multi-modal edge AI,''
\emph{ACM Conference on Mobile Systems, Applications, and Services (MobiSys)}, 2024.

\bibitem{hu2021lora}
E.~J. Hu, Y.~Shen, P.~Wallis, Z.~Allen-Zhu, Y.~Li, S.~Wang, L.~Wang, and W.~Chen,
``LoRA: Low-rank adaptation of large language models,''
\emph{International Conference on Learning Representations (ICLR)}, 2022.

\end{thebibliography}

\end{document}
